\documentclass[11pt, addpoints]{exam}
\usepackage{medipol}
% For an exam, single spacing is most appropriate
\singlespacing
% \onehalfspacing
% \doublespacing
% For an exam, we generally want to turn off paragraph indentation
\parindent 0ex
\usehorizontalhalf
\bracketedpoints
\qformat{Question \thequestion\dotfill \emph{\totalpoints\ points}}
\settabletotalpoints{100}

%%%%%%%%%%%%%%%%%%%%%%%% Put here your packages and commands %%%%%%%%%%%%%%

\usepackage[most]{tcolorbox}
\tcbuselibrary{theorems}
\tcbset{colback=red!1!white,colframe=red!5!black,width=4cm,height=1cm, box align=center}

%%%%%%%%%%%%%%%%%%%%%%%%% Exam Data %%%%%%%%%%%%%%%%%

\coursename{Course name or Course code} % Course name or code
\school{Engineering and Natural Sciences}% Faculty or School
\department{Computer Engineering}% Department
\semester{G\"{u}z / Fall}% Semester
\year{2023--2024}% Academic year
\examnum{1\textsuperscript{st} Midterm}% Midterm or Final exam
\datetime{Date and Time}% Date and time of the exam
\timelimit{Duration of exam}% Duration of the exam

%%%%%%%%%%%%%%%%%%%%%%%%%% Defining your categories %%%%%%
% These are used to introduce the learning outcomes of  the exam

\newcategory{1.1}{Statement of Learning Outcome 1.1}
\newcategory{1.2}{Statement of Learning Outcome 1.2}
\newcategory{1.3}{Statement of Learning Outcome 1.3}
\newcategory{2.1}{Statement of Learning Outcome 2.1}
\newcategory{2.2}{Statement of Learning Outcome 2.2}
\newcategory{2.3}{Statement of Learning Outcome 2.3}
\newcategory{3}{Statement of Learning Outcome 3}
\newcategory{4}{Statement of Learning Outcome 4}
\newcategory{5.1}{Statement of Learning Outcome 5.1}
\newcategory{5.2}{Statement of Learning Outcome 5.2 (very long very long very long very long very long \\ & 
very long very long very long very long very long very long very long very long \\ & very long very long very long)}
\newcategory{5.3}{Statement of Learning Outcome 5.3}


\begin{document}

\coverpage%this is used to print Medipol's header and footer on the current page.

\begin{itemize}
\item  Fill the exam data in the preamble:
\begin{verbatim}
\coursename{THE COURSE NAME OR CODE}
\school{Engineering and Natural Sciences}
\department{THE DEPARTMENT}
\semester{THE SEMESTER} % For example: Güz / Fall
\year{THE ACADEMIC YEAR} % For example: 2023 -- 2024
\examnum{EXAM TYPE}% Midterm or Final exam
\datetime{DATE AND TIME OF THE EXAM}% For example: November 26,2023, 09:00 A.M.
\timelimit{DURATION OF THE EXAM}
\end{verbatim}
\item The command \verb|\newpage| produces a new page similar to page 3 in the present file (without Medipol's header and footer).
\item The commands \verb|\newpage\coverpage| produce a new page similar to page 2 in the present file (with Medipol's header and footer).
\item In the preamble, define each one of your learning outcomes as follows:
\begin{verbatim}
\newcategory{OUTCOME NO.}{STATEMENT OF THE OUTCOME} 
% For example: \newcategory{2.3}{Differentiate power functions.} 
\end{verbatim}
\item To assign a learning outcome to a question (part, or subpart) use \verb|\categorize{OUTCOME NO.}|
\item To print the grading table use \verb|\gradetable[h]|. If you use \verb|\gradetable[v]|, then the learning outcomes will be ignored from the grading table.
\item To print the table of learning outcomes use \verb|\categorytable|
\item If the statement of one of the outcomes is longer than the line width and it makes the table of learning outcomes ugly, you can divide it into several lines using \verb|\\ &|. See outcome 5.2 in the table on page 7 in the present file.
\end{itemize}

%%%%%%%%%%%%%%%%% Ex: Page 2 %%%%%%%%%%%%5

\newpage
\coverpage

\begin{center}
{\Large{\color{red}{Ex: Cover page with Medipol's header and footer.}}}\\
\end{center}



This exam contains \numpages\ pages (including this cover page) and
\numquestions\ problems.  Make sure you attempt all the questions.  Make sure to type your initials on every page of answer sheet.\\

%You may \textit{not} use your books, notes, or any calculator on this exam
This exam is a closed book exam. Neither calculators nor phones are allowed. You are required to show your work on each problem on this exam.  The following rules apply:

\begin{itemize}
\item Provide your answers after each question for the space provided.
\item \textbf{Mysterious or unsupported answers will not receive full
credit}.  A correct answer, \textbf{unsupported by calculations}, explanation,
or algebraic work will receive no credit; an incorrect answer supported
by \textbf{substantially correct calculations} and explanations might still receive
partial credit.
\item Make sure you pay attention to bold text.
\item Do not write in the table below.
\end{itemize}

\vspace{0.5cm}

\begin{center}
\gradetable[h]
\end{center}

\vspace{0.5cm}

\begin{flushright}
\begin{tabular}{@{}c@{}}
With best wishes,\\

Prof.\ Dr.\ first instructor\\
Dr.\ second instructor\\
Dr.\ third instructor\\
\end{tabular}
\end{flushright}    

%%%%%%%%%%%%%%%%% Ex: Page 3 %%%%%%%%%%%%5

\newpage

\begin{center}
{\Large{\color{red}{Ex: A page without Medipol's header and footer, and a question with parts.}}}\\
\end{center}

\begin{questions}
\addpoints 
\question  Example of a question with parts:
\begin{parts}
\part[10]\categorize{2.1} Part with learning outcome 2.1
\vfill
\part[10] \categorize{1.1} Part with learning outcome 1.1
\vfill
\part[5]\categorize{5.1}\categorize{2.2} Part with learning outcomes 5.1 and 2.2
\vfill
\end{parts}

%%%%%%%%%%%%%%%%% Ex: Page 4 %%%%%%%%%%%%5

\newpage

\begin{center}
{\Large{\color{red}{Ex: A page without Medipol's header and footer, and a question with parts and subparts.}}}\\
\end{center}

\addpoints 
\question  Example of a question with subparts:
\begin{parts}
\part
\begin{subparts}
 \subpart[7] \categorize{5.1} Subpart with learning outcome 5.1
\vfill
\subpart[8] \categorize{5.2} Subpart with learning outcome 5.2
\vfill
\end{subparts}
\part[10]\categorize{5.3} Part with learning outcome 5.3
\vfill
\end{parts}

%%%%%%%%%%%%%%%%% Ex: Page 5 %%%%%%%%%%%%5

\newpage
\coverpage

\begin{center}
{\Large{\color{red}{Ex: A page with Medipol's header and footer, and a question without parts.}}}\\
\end{center}


\addpoints 
\question[25]\categorize{1.2}\categorize{1.3}\categorize{2.3} Question with learning outcomes 1.2, 1.3 and 2.3

\vfill


%%%%%%%%%%%%%%%%% Ex: Page 6 %%%%%%%%%%%%5

\newpage
\coverpage

\begin{center}
{\Large{\color{red}{Ex: A page with Medipol's header and footer, and a question with parts and subparts.}}}\\
\end{center}


\addpoints 

\question  Example of a question with subparts:
\begin{parts}
\part
\begin{subparts}
 \subpart[5] \categorize{1.2} Subpart with learning outcome 1.2
\vfill
\subpart[10] \categorize{2.3} Subpart with learning outcome 2.3
\vfill
\end{subparts}
\part[10]\categorize{1.3} Part with learning outcome 1.3
\vfill
\end{parts}

\end{questions}

%%%%%%%%%%%%%%%%% Ex: Page 7 %%%%%%%%%%%%5

\newpage

\begin{center}
{\Large{\color{red}{Ex: Extra page without Medipol's header and footer. You can use this for formula sheets and/or list of outcomes measured in the exam.}}}\\
\end{center}


\begin{center}
\categorytable    
\end{center}



\end{document}